\item \textbf{Energía y radiación residual de uranio empobrecido en obuses.}
El uranio empobrecido ($^{238}\text{U}$) se utiliza en obuses de artillería antiblindaje debido a su extrema densidad de masa ($\rho = 19.1 \times 10^3\text{ kg/m}^3$).
\begin{enumerate}
    \item Encuentre la arista de un cubo macizo de uranio de exactamente $70.0\text{ kg}$ de masa.
    \item El isótopo $^{238}_{92}\text{U}$ decae lentamente ($T_{1/2} = 4.47 \times 10^9\text{ años}$) iniciando una serie rápida de 14 etapas sucesivas que resultan en la reacción neta global:
    $$ ^{238}_{92}\text{U} \to 8(^{4}_{2}\text{He}) + 6(e^-) + ^{206}_{82}\text{Pb} + 6(\bar{\nu}) + Q_{\text{neta}} $$
    Calcule la energía neta de desintegración $Q_{\text{neta}}$ liberada.
    \item Demuestre que la potencia térmica de salida de una muestra con actividad $R$ y energía de desintegración $Q$ es $P = QR$.
    \item Calcule la potencia radiactiva de salida (en $\text{J/año}$) para el obús de $70.0\text{ kg}$, asumiendo que ha alcanzado el equilibrio secular con sus isótopos hija.
    \item \textbf{¿Qué pasaría si?} Un soldado trabaja en un arsenal expuesto a la radiación de estos obuses, con un límite anual de dosis equivalente de $5.00\text{ rem/año}$. Determine la tasa máxima a la que puede absorber energía de radiación en $\text{J/año}$ si su masa es de $70.0\text{ kg}$ y el factor promedio de efectividad biológica relativa (RBE) de la radiación es $1.10$.
\end{enumerate}